\documentclass[11pt,a4paper]{article}
\usepackage[T1]{fontenc}
\usepackage[utf8]{inputenc}
\usepackage{lmodern,amsmath,amssymb,amsthm,mathtools,microtype,mathrsfs}
\usepackage[margin=25mm,top=23mm,bottom=25mm]{geometry}
\usepackage{booktabs,array,longtable,enumitem,xcolor,fancyhdr}
\usepackage[colorlinks=true,linkcolor=blue!45!black,urlcolor=blue!45!black,citecolor=blue!45!black]{hyperref}
\hypersetup{pdftitle={Reading the Alpoge--Claude RH and Dirichlet-GRH proof: working version 0.35},pdfauthor={ChatGPT},pdfsubject={Partial human-readable reconstruction with explicit arithmetic inputs},pdfkeywords={Riemann hypothesis, Lean, theta lifts, prime sums, Arakelov slopes}}
\setlength{\parindent}{0pt}
\setlength{\parskip}{5pt plus 1pt}
\setlist{topsep=3pt,itemsep=3pt,leftmargin=*}
\setlength{\headheight}{14pt}
\pagestyle{fancy}\fancyhf{}
\fancyhead[L]{\small RH and Dirichlet GRH: a reconstruction}
\fancyhead[R]{\small Working version 0.35}
\fancyfoot[C]{\thepage}
\renewcommand{\headrulewidth}{0.3pt}
\numberwithin{equation}{section}
\newtheorem{theorem}{Theorem}[section]
\newtheorem{proposition}[theorem]{Proposition}
\newtheorem{lemma}[theorem]{Lemma}
\newtheorem{corollary}[theorem]{Corollary}
\theoremstyle{remark}\newtheorem{remark}[theorem]{Remark}
\newcommand{\Rea}{\operatorname{Re}}
\newcommand{\HS}{\mathrm{HS}}
\newcommand{\tr}{\operatorname{tr}}
\newcommand{\codim}{\operatorname{codim}}
\newcommand{\covol}{\operatorname{covol}}
\newcommand{\OK}{\mathcal O_K}
\newcommand{\src}[2]{\href{https://github.com/CocoSgt/riemann-hypothesis-lean/blob/7cbf7a0c32fca1badf05931879a37c3f8a2b03b4/src/#1}{\texttt{#2}}}
\newcommand{\status}[1]{\textbf{Status:} #1}
\begin{document}
\section{The original regular quotient, smooth intertwiners and cuspidal closure}
\label{sec:regular-quotient}

This chapter supplies analytic connections used after HCU18.
Its finite-volume input is exactly the source's \(\mathrm{FinPair}\);
it does not construct that input.  Initial boundedness of the
holomorphic automorphic function is still a separate obligation.
All source statements refer to commit
\texttt{7cbf7a0c32fca1badf05931879a37c3f8a2b03b4}.

Fix the original datum \(D\) and a nondegenerate Hermitian space \(V\).
Write \(G=\mathrm{GA}(D,V)\), \(\Gamma=\mathrm{GQ}(V)\),
\(Q=\mathrm{AutQ}(V)=\Gamma\backslash G\), and
\(\pi=\mathrm{autMk}(V)\).  A supplied finite pair consists of a
left Haar measure \(\mu\) which is also right invariant and a null
measurable fundamental domain \(\mathcal F\) for the left action of
\(\Gamma\), with \(\mu(\mathcal F)<\infty\).  Throughout,
\[
 \nu=(\mu|_{\mathcal F})_*\pi,\qquad
 \mathcal H=L^2(Q,\nu).
 \tag{RQC.1}
\]
This is the literal \(\mathrm{nuQ}\), including its scale.
The condition \(\mathrm{Repr}(F,f)\) is equality
\(F(\pi x)=f(x)\) for \(\mu\)-almost every \(x\) in all of \(G\),
not only on the fundamental domain.

% BEGIN RQC1
\paragraph{RQC1: topology of the original adelic group.}
\label{claim:RQC1}
The original \(G\) is a locally compact, second-countable Hausdorff
topological group with its stated Borel structure.  Its rational
subgroup \(\Gamma\) is countable, discrete and closed.

\emph{Proof.}
There are countably many finite places of a number field.  At a finite
place the ring of integers is compact: completeness identifies it
with the inverse limit of the finite quotients by powers of a
uniformizer.  The local field is the countable union of dilates of
this compact open set.  A countable dense copy of the number field,
together with valuation balls, gives a countable base.  The
restricted product is a countable union, over finite sets of places,
of open products of finitely many local fields with the remaining
compact integer rings.  Each such product is locally compact and
second countable.  Adding the finitely many archimedean factors
gives these properties for \(\mathbb A_F\).

A finite \(F\)-basis of the original coefficient field \(E\)
identifies \(\mathbb A_F\otimes_F E\), with its source module
topology, with finitely many copies of \(\mathbb A_F\).  For the
original nonsingular Gram matrix \(H\), the matrix realization is
\[
 G=\{g:g^*Hg=H\},\qquad g^{-1}=H^{-1}g^*H.
 \tag{RQC.2}
\]
The equation implies invertibility by taking determinants, and the
displayed inverse is continuous.  Thus this closed matrix subspace
has precisely the topology inherited from the group of units used
by the source.  Closed subspaces preserve local compactness and
second countability; multiplication and inversion are continuous.

For discreteness, a rational \(a\in F\) integral at every finite
place and of absolute value less than one at every archimedean
embedding must be zero: otherwise its nonzero integral norm has
absolute value less than one.  These conditions give an adelic
neighborhood of zero meeting \(F\) only at zero.  The finite basis
gives the analogous neighborhood for \(E\).  Apply it to every
entry of \(g-1\) to obtain an identity neighborhood meeting
\(\Gamma\) only at \(1\).  Rational matrices form a countable set.
A discrete subgroup of a Hausdorff topological group is closed:
choose \(U\) with \(U^{-1}U\cap\Gamma=\{1\}\); each translate of
\(U\) contains at most one rational element, so any closure point
in that translate equals that element by the Hausdorff property.
Finally the quotient map is continuous and open, and right
multiplication descends to continuous maps on its quotient Borel
space.  No finite-volume assertion has been used.
% END RQC1

% BEGIN RQC2
\paragraph{RQC2: unfolding and independence of the fundamental domain.}
\label{claim:RQC2}
For nonnegative measurable \(h\) on \(G\),
\[
 \int_G h\,d\mu
 =\sum_{\gamma\in\Gamma}\int_{\mathcal F}h(\gamma x)\,d\mu(x).
 \tag{RQC.3}
\]
A left \(\Gamma\)-invariant function has the same integral on any
two fundamental domains, first as an extended nonnegative integral
and then as a complex integral whenever it is integrable on one.

\emph{Proof.}
Countability makes the exceptional set on which the translates
\(\gamma\mathcal F\) fail to partition \(G\) a single null set.
Countable additivity and left Haar invariance give (RQC.3).
For a second domain \(\mathcal E\), partition it into
\(\mathcal E\cap\gamma\mathcal F\), translate each piece by
\(\gamma^{-1}\), and use invariance of \(h\).
The resulting pieces \(\gamma^{-1}\mathcal E\cap\mathcal F\)
partition \(\mathcal F\) almost everywhere.  Tonelli gives the
nonnegative equality.  Applying it to \(|h|\), then to the positive
and negative parts of the real and imaginary parts, gives the
complex assertion.  Null measurable domains are treated by
measurable representatives and their countably many translated null
differences.
% END RQC2

% BEGIN RQC3
\paragraph{RQC3: the global almost-everywhere transfer.}
\label{claim:RQC3}
An equality holding \(\nu\)-almost everywhere on \(Q\) holds after
composition with \(\pi\), \(\mu\)-almost everywhere on all of \(G\).

\emph{Proof.}
For a null exceptional set \(A\subset Q\), (RQC.1) gives
\(\mu(\mathcal F\cap\pi^{-1}A)=0\).
Its lift is left \(\Gamma\)-invariant, so its intersection with
every \(\gamma\mathcal F\) is null.  These countably many translates
cover almost all of \(G\).  A measurable null cover handles
completed-measure exceptional sets.  This proves the source's
global almost-everywhere convention.
% END RQC3

% BEGIN RQC4
\paragraph{RQC4: right translation preserves the actual quotient measure.}
\label{claim:RQC4}
Set \(\tau_g(\pi x)=\pi(xg)\).  For every \(g\in G\), this is a
measurable measure-preserving bijection of \((Q,\nu)\), and
\(\tau_h\circ\tau_g=\tau_{gh}\).

\emph{Proof.}
Left rational equivalence commutes with right multiplication.
The inverse is \(\tau_{g^{-1}}\), and RQC1 gives continuity.
For nonnegative measurable \(a:Q\to[0,\infty]\), right Haar
invariance gives
\[
 \int_Q a\circ\tau_g\,d\nu
 =\int_{\mathcal F}a(\pi(xg))\,d\mu(x)
 =\int_{\mathcal Fg}a(\pi x)\,d\mu(x).
 \tag{RQC.4}
\]
The right-translated domain is another left rational fundamental
domain.  RQC2 makes the last expression \(\int_Q a\,d\nu\).
The composition law follows by evaluating at \(\pi x\).
No Haar normalization or original measure has been changed.
% END RQC4

% BEGIN RQC5
\paragraph{RQC5: the comparison estimate and strong continuity.}
\label{claim:RQC5}
For measurable \(A\subset Q\), put \(B=\mathcal F\cap\pi^{-1}A\).
Then
\[
 \nu(\tau_g^{-1}A\mathbin{\triangle}\tau_h^{-1}A)
 \le\mu(Bg^{-1}\mathbin{\triangle}Bh^{-1}).
 \tag{RQC.5}
\]
Right pullback therefore acts strongly continuously on \(L^2(Q,\nu)\).

\emph{Proof.}
For measurable \(S\subset G\), unfolding gives
\[
 \mu(\mathcal F\cap\Gamma S)
 \le\sum_{\gamma\in\Gamma}\mu(\mathcal F\cap\gamma S)
 =\mu(S).
 \tag{RQC.6}
\]
Apart from a null set, if \(\pi(xg)\in A\) and
\(\pi(xh)\notin A\), move \(xg\) into \(\mathcal F\) by a
rational element \(\gamma\).  Then
\(\gamma x\in Bg^{-1}\setminus Bh^{-1}\).  Reverse the roles for
the other case.  Thus the lift of the left symmetric difference
lies almost everywhere in the rational saturation of the right
one.  Applying (RQC.6) proves (RQC.5).  The exceptional sets stay
null under the two fixed right translations by right invariance.

For a finite-measure set \(B\) in this locally compact group,
\(\mu(Bu\triangle B)\to0\) as \(u\to1\).
Indeed Haar regularity gives compact \(K\subset B\) and open
\(O\supset K\) with both \(\mu(B\setminus K)\) and
\(\mu(O\setminus K)\) arbitrarily small, and \(\mu(O)<\infty\).
Compactness gives \(Ku\subset O\) for \(u\) near \(1\).  Hence
\[
 \mu(Ku\triangle K)\le2\mu(O\setminus K),\qquad
 \mu(Bu\triangle B)\le
 2\mu(B\setminus K)+2\mu(O\setminus K).
 \tag{RQC.7}
\]
The argument also applies to null measurable sets after replacement
by measurable representatives.  In (RQC.5), \(B\) has finite
measure.  The square of the \(L^2\) distance between the two
pulled-back indicators is its left side, so indicator orbits and
then simple-function orbits are continuous.  Simple functions are
dense in \(L^2\), by truncation and step approximation of real and
imaginary parts.  Pullback is an isometry by RQC4; replacing a
vector by a simple function changes an orbit difference by at
most twice the approximation error, uniformly in the translating
element.  This proves continuity for all vectors.
% END RQC5

% BEGIN RQC6
\paragraph{RQC6: the full original regular-representation producer.}
\label{claim:RQC6}
On the same space \(\mathcal H\), there is a monoid homomorphism
into continuous complex-linear endomorphisms,
\[
 \rho:G\longrightarrow\mathcal L_{\mathbb C}(\mathcal H),
 \qquad \rho(g)F=F\circ\tau_g,
 \tag{RQC.8}
\]
satisfying every field of \(\mathrm{IsReg}\): the value identity
holds globally \(\mu\)-almost everywhere, each operator is
isometric, and every orbit is continuous.

\emph{Proof.}
RQC4 makes pullback well defined on equivalence classes, complex
linear and norm preserving.  On representatives,
\(\rho(g)\rho(h)F(\pi x)=F(\pi(xgh))\), so its homomorphism
order is exactly the original one.  The identity and inverse
follow similarly.  RQC3 promotes the quotient almost-everywhere
value identity to all of \(G\), yielding the literal
\texttt{val} field.  RQC4 gives \texttt{iso}, and RQC5 gives
\texttt{cont}.  This is \texttt{QuotRegular.exists\_isReg} from the
supplied \(\mathrm{FinPair}\), with no further analytic input.
Moreover any supplied \(\mathrm{IsReg}\) on this quotient equals
this representation: its global value law restricts to
\(\mathcal F\), hence gives the same \(\nu\)-almost-everywhere
pullback for each \(g,F\).  Thus the arguments below also apply
to any original witness of \(\mathrm{IsReg}\).
% END RQC6

\subsection{Original smooth vectors and all derivative words}

For this subsection \(H\) is a nonsingular Hermitian matrix, and
\(\sigma\) is a strongly continuous unitary representation of
\(U(H)\) on a complex Hilbert space \(\mathcal K\).
The original \(\mathrm{SmoothVec}\) condition says that for every
finite list \(X_1,\ldots,X_r\in
\mathfrak u(H)=\{X:X^*H+HX=0\}\), the map
\[
 (t_1,\ldots,t_r)\longmapsto
 \sigma(e^{t_1X_1}\cdots e^{t_rX_r})v
 \tag{RQC.9}
\]
is real smooth.  The original derivative is
\(d\sigma(X)v=\frac{d}{dt}|_0\sigma(e^{tX})v\).

% BEGIN RQC7
\paragraph{RQC7: the full smooth subspace and derivative stability.}
\label{claim:RQC7}
Smooth vectors form a complex subspace, invariant under
\(\sigma(g)\), every real derived operator, and every complexified
derived operator.  Thus the original span \(M.\mathrm{sm}\)
contains only smooth vectors.

\emph{Proof.}
Finite sums and complex scalar multiples commute with the bounded
linear operators \(\sigma(g)\) and with derivatives in (RQC.9).
Group invariance follows from
\[
 \sigma(e^{t_1X_1}\cdots e^{t_rX_r})\sigma(g)v
 =\sigma(g)\sigma(e^{t_1g^{-1}X_1g}\cdots
                       e^{t_rg^{-1}X_rg})v.
 \tag{RQC.10}
\]
Conjugation by \(g\in U(H)\) preserves \(\mathfrak u(H)\).
Also \(e^{tX}\in U(H)\), since differentiating
\(e^{tX^*}He^{tX}\) gives zero.

Append \(e^{sX}\) to (RQC.9).  The resulting function of
\((t,s)\) is smooth, and its partial derivative at \(s=0\) is
\(\sigma(e^{t_1X_1}\cdots e^{t_rX_r})d\sigma(X)v\):
the preceding representation operator is bounded linear.
Partial derivatives of smooth Hilbert-valued functions are smooth,
proving derivative stability without an interchange of unbounded
operators.

Set \(c_H(Z)=-H^{-1}Z^*H\).  This is a conjugate-linear involution
with fixed real subspace \(\mathfrak u(H)\), so
\[
 X=\tfrac12(Z+c_HZ),\quad
 Y=\tfrac1{2i}(Z-c_HZ),\quad Z=X+iY
 \tag{RQC.11}
\]
has \(X,Y\in\mathfrak u(H)\).
The source complexified derivative is \(d\sigma(X)+i\,d\sigma(Y)\).
Its stability follows from what was just proved.  These statements
concern the entire original smooth subspace, not a selected dense
subspace.
% END RQC7

% BEGIN RQC8
\paragraph{RQC8: complete derived Lie algebra and adjoint identities.}
\label{claim:RQC8}
On smooth vectors, \(Z\mapsto d\sigma_{\mathbb C}(Z)\) is complex
linear with complex-linear endomorphism values, and
\[
\begin{split}
 [d\sigma_{\mathbb C}(Z),d\sigma_{\mathbb C}(W)]
   &=d\sigma_{\mathbb C}([Z,W]),\\
 \sigma(g)d\sigma_{\mathbb C}(Z)
   &=d\sigma_{\mathbb C}(gZg^{-1})\sigma(g),\\
 \langle d\sigma_{\mathbb C}(Z)v,w\rangle
   &=\langle v,d\sigma_{\mathbb C}(H^{-1}Z^*H)w\rangle.
\end{split}
 \tag{RQC.12}
\]
The inner product is conjugate linear in its first argument.
On a real Lie direction the complexified derivative is the original
real derivative.  These are all clauses of the source's
\texttt{complexifiedDerivedAction\_isLieAction}.

\emph{Proof.}
We give the chart justification for linearity in the direction.
Near zero in \(\mathfrak u(H)\), the Cayley map
\(X\mapsto(1+X)(1-X)^{-1}\) has values in \(U(H)\).
Its inverse near \(1\) is \((g-1)(g+1)^{-1}\), and its derivative
at zero is \(2X\).  Choose a real basis \(X_1,\ldots,X_d\).
In this chart, \(t\mapsto\prod_j e^{t_jX_j}\) has derivative
\(t\mapsto\frac12\sum_jt_jX_j\), an isomorphism.
The real inverse function theorem gives local smooth coordinates
of this product form.  Therefore (RQC.9) makes
\(g\mapsto\sigma(g)v\) smooth in the group chart.  Its differential
at \(1\) is real linear and takes \(X\) to \(d\sigma(X)v\).
Vector sums give complex linearity in \(v\); (RQC.11) then gives
complex linearity in \(Z\) and the real restriction.

Apply left-invariant vector fields \(L_X(g)=gX\) and \(L_Y(g)=gY\)
to this orbit map.  Its first derivative is
\(L_X(\sigma(\cdot)v)(g)=\sigma(g)d\sigma(X)v\).
In the difference of second derivatives, the symmetric second
differential cancels, and differentiating the fields gives
\(g(XY-YX)\).  At \(1\) this says
\([d\sigma(X),d\sigma(Y)]v=d\sigma([X,Y])v\).
RQC7 supplies the required domains.  Expanding (RQC.11) gives
the complex bracket identity.

Differentiating
\(\sigma(g)\sigma(e^{tX})=\sigma(e^{tgXg^{-1}})\sigma(g)\)
gives real covariance, then complex covariance.
Differentiating preservation of
\(\langle\sigma(e^{tX})v,\sigma(e^{tX})w\rangle\) gives
\(\langle d\sigma(X)v,w\rangle=-\langle v,d\sigma(X)w\rangle\).
For \(Z=X+iY\), the resulting adjoint direction is
\(-X+iY=-c_H(Z)=H^{-1}Z^*H\), proving the final identity.
All calculations are on smooth vectors; no boundedness of a
derived operator is asserted.
% END RQC8

% BEGIN RQC9
\paragraph{RQC9: the literal equivariant multilinear smoothness producer.}
\label{claim:RQC9}
Fix the original \(\theta\) satisfying \(\mathrm{IsInfFam}\),
model \(M\), and a monoid homomorphism \(\rho\) into continuous
complex-linear operators on any complex inner-product space
\(\mathcal K\).  Let
\(\Phi:\prod_\tau M.H_\tau\to\mathcal K\) be continuous complex
multilinear, with the original separate \(\mathrm{atTau}\)
equivariance.  For \(y_\tau\in M.\mathrm{sm}_\tau\),
\(\Phi(y)\) is an original \(\mathrm{IsSmoothVector}'\), and
\[
 dL^2(\tau,X)\Phi(y)
 =\Phi\bigl(y[\tau\leftarrow d(M.\mathrm{rep}_\tau)(X)y_\tau]\bigr)
 \quad(X\in\mathrm{lieAt}(V,\tau)).
 \tag{RQC.13}
\]

\emph{Proof.}
The source action is
\(M.\mathrm{rep}_\tau(g)=
 M.\sigma_\tau(M.\mathrm{fr}_\tau gM.\mathrm{fr}_\tau^{-1})\).
The original frame equation identifies the evaluated Gram form
with the signed model form.  It identifies their unitary groups,
Lie algebras, products and exponentials, so RQC7 applies to this
action and the entire original smooth span.
For an ordered list \((\tau_j,X_j)\), HCU14 identifies the actual
\(\mathrm{oneP}\) with the same one-place lift of \(e^{t_jX_j}\).
Repeated separate equivariance gives
\[
 \rho\!\left(\prod_j\mathrm{oneP}(\tau_j,X_j,t_j)\right)\Phi(y)
 =\Phi\left(\left(
 M.\mathrm{rep}_\tau\!\left(
       \prod_{j:\tau_j=\tau}e^{t_jX_j}\right)y_\tau
                         \right)_\tau\right).
 \tag{RQC.14}
\]
Both products retain their original order.  Each coordinate on the
right is smooth by (RQC.9), inserting zero directions at unused
positions.  A continuous multilinear map is smooth: expand its
values after adding increments in each variable to obtain its
finite Taylor expansion and continuous multilinear derivatives.
Thus (RQC.14) is precisely the all-word smoothness predicate.
For a single letter, differentiate through the bounded linear
partial map of \(\Phi\) to obtain (RQC.13), with the exact
original \(dL^2\).  This is the whole endpoint
\texttt{SmoothMultilinear.smooth\_map}.
% END RQC9

% BEGIN RQC10
\paragraph{RQC10: every Hilbert derivative word with its original order.}
\label{claim:RQC10}
For a list \(l\) of real Lie pairs \((\tau,X)\), let
\(l_\tau=\mathrm{lettersAt}(\tau,l)\) be its ordered filter at
\(\tau\).  Then
\[
 l.\mathrm{foldr}(dL^2)\,\Phi(y)=
 \Phi\left(\left(
 l_\tau.\mathrm{foldr}\bigl(d(M.\mathrm{rep}_\tau)\bigr)y_\tau
                         \right)_\tau\right),
 \tag{RQC.15}
\]
and all resulting coordinate vectors remain in the original smooth
subspaces.

\emph{Proof.}
For the empty word this is the identity.  Apply the induction
hypothesis to the tail, then (RQC.13) to the leading pair
\((\tau,X)\).  RQC7 keeps the new coordinate smooth.
Filtering puts \(X\) at the beginning of that place's word and
leaves the other filters unchanged.  This is exactly the
right-fold order, without commuting directions at the same place.
It also proves the source's \texttt{word\_mem} and
\texttt{word\_map}.  The finite eigenlaw in HCU18 holds for every
Hilbert tuple, so it holds for each of these derived tuples with
the same finite subgroup and character.  The level is not chosen
again as the derivative order grows.
% END RQC10

\subsection{Compact averaging and the original cusp-closure argument}

% BEGIN RQC11
\paragraph{RQC11: compact orbit multiplicity and a local integral bound.}
\label{claim:RQC11}
For compact \(K\subset G\), put
\(N_K=\#(\Gamma\cap KK^{-1})<\infty\).
Every \(x\in G\) lies in at most \(N_K\) of the sets \(\gamma K\).
For measurable left rational invariant \(h\),
\[
 \int_K|h|^p\,d\mu\le N_K\int_{\mathcal F}|h|^p\,d\mu
 \quad(p=1,2).
 \tag{RQC.16}
\]
If \(h\) represents \(H_0\in\mathcal H\), then
\[
 \int_K|h|\,d\mu
 \le N_K\mu(\mathcal F)^{1/2}\|H_0\|_2.
 \tag{RQC.17}
\]

\emph{Proof.}
RQC1 makes \(\Gamma\cap KK^{-1}\) finite.
If \(\gamma_0^{-1}x\in K\), every other
\(\gamma^{-1}x\in K\) determines a distinct
\(\gamma^{-1}\gamma_0\in\Gamma\cap KK^{-1}\).
Unfold \(1_K|h|^p\) to get
\[
 \int_K|h|^p\,d\mu
 =\int_{\mathcal F}|h(x)|^p
       \sum_{\gamma\in\Gamma}1_K(\gamma x)\,d\mu(x),
\]
and use this multiplicity bound.  The representative identity
gives \(\int_{\mathcal F}|h|^2=\|H_0\|_2^2\).
Cauchy--Schwarz on the finite domain gives (RQC.17).
No pointwise bound on an arbitrary \(L^2\) class is claimed.
% END RQC11

% BEGIN RQC12
\paragraph{RQC12: compact averaging is locally controlled in \(L^1\).}
\label{claim:RQC12}
Let \(m\) be a finite Borel measure on \(G\), concentrated on
compact \(K\).  For continuous \(h:G\to\mathbb C\), the function
\(A_mh(g)=\int h(zg)\,dm(z)\) is continuous.  For compact
\(\Omega\subset G\),
\[
 \int_\Omega|A_mh(g)|\,d\mu(g)
 \le m(G)\int_{K\Omega}|h(x)|\,d\mu(x).
 \tag{RQC.18}
\]
If \(h\) is rational invariant and represents \(H_0\), the last
quantity is at most
\(m(G)N_{K\Omega}\mu(\mathcal F)^{1/2}\|H_0\|_2\).

\emph{Proof.}
On \(K\) times a compact neighborhood of a point, continuity
gives a uniform bound and local uniform convergence in \(g\);
integration against finite \(m\) therefore preserves continuity.
The triangle inequality, Tonelli and left Haar invariance give
\[
 \int_\Omega|A_mh|\,d\mu
 \le\int_K\int_\Omega|h(zg)|\,d\mu(g)\,dm(z)
 =\int_K\int_{z\Omega}|h(x)|\,d\mu(x)\,dm(z).
\]
Use \(z\Omega\subset K\Omega\), then RQC11.
These are ordinary integrals of continuous functions on compact
sets.
% END RQC12

% BEGIN RQC13
\paragraph{RQC13: vanishing compact averages survive closed orbit spans.}
\label{claim:RQC13}
Let continuous rational invariant \(\phi\) represent \(F_0\),
with \(A_m\phi(g)=0\) for every \(g\).
If \(F\) lies in the closed complex span of
\(\{\rho(s)F_0:s\in G\}\), every continuous rational invariant
representative \(\psi\) of \(F\) satisfies \(A_m\psi=0\) everywhere.

\emph{Proof.}
Every finite orbit combination \(H_j\) has a continuous rational
invariant representative \(h_j\), formed from the same combination
of right translates of \(\phi\).  RQC6 gives its global
\(\mathrm{Repr}\) identity, and
\(A_m(\phi(\cdot s))(g)=A_m\phi(gs)=0\).
Choose \(H_j\to F\) in \(L^2\).  The function \(\psi-h_j\)
represents \(F-H_j\).  RQC12 yields
\[
 \int_\Omega|A_m\psi|\,d\mu
 \le m(G)N_{K\Omega}\mu(\mathcal F)^{1/2}\|F-H_j\|_2
 \longrightarrow0
 \tag{RQC.19}
\]
for every compact \(\Omega\).  A nonzero value of the continuous
function \(A_m\psi\) would give a relatively compact open
neighborhood on which its absolute value has a positive lower
bound.  Haar measure is positive on that neighborhood, contradicting
(RQC.19) on its compact closure.

The proof controls averages by local \(L^1\) estimates.  It never
restricts an arbitrary ambient \(L^2\) equivalence class to a
measure-zero unipotent subgroup.  The assumed continuous
representative is an explicit part of the statement.
% END RQC13

% BEGIN RQC14
\paragraph{RQC14: a compactly contained domain for the actual \(ZI\).}
\label{claim:RQC14}
For every original nonzero isotropic rational vector \(\xi\),
\(ZI(V,\xi)\) has a Borel rational fundamental domain contained
in a compact subset, valid for every Haar measure.  On each fixed
right translate, a rational-invariant continuous function has the
same absolutely convergent integral over this domain as over any
supplied rational fundamental domain.

\emph{Proof.}
CTC1, applied to the original nondegenerate target form, gives
\[
 \mathbb A_F\xrightarrow{\ \sim\ }ZI(V,\xi),\qquad
 t\longmapsto1+\delta tN_\xi,\quad
 N_\xi=\xi\xi^*H,\quad N_\xi^2=0,\quad
 \bar\delta=-\delta\ne0.
 \tag{RQC.20}
\]
It identifies rational points with \(F\) and proves closedness and
cocompactness; no definiteness of the target was used.  There is
therefore a compact set \(C\) meeting every rational orbit.

Choose an injective Borel \(\kappa:C\to\mathbb R\) by encoding
membership in a countable separating base with ternary digits
zero and two.  Each orbit meets \(C\) in a finite set, since the
rational subgroup is closed discrete.  Select in that set the
point of least \(\kappa\).  Explicitly keep \(x\in C\) for which
\(\kappa(x)\le\kappa(\gamma x)\) whenever \(\gamma x\in C\).
Countably many Borel conditions define a Borel set
\(\mathcal E\subset C\), meeting each orbit exactly once.
Its translates partition the group; thus it is a fundamental
domain for every Haar measure \(\eta\), has compact closure,
and \(\eta(\mathcal E)<\infty\).

For fixed \(g\), the function \(z\mapsto h(zg)\) is rational
invariant and bounded on \(C\).  The partition argument RQC2,
now within \(ZI(V,\xi)\), identifies its integral and its absolute
value integral on \(\mathcal E\) with those on any supplied
\(\eta\)-fundamental domain.  Absolute convergence follows.
Neither the supplied measure nor its normalization is replaced.
% END RQC14

% BEGIN RQC15
\paragraph{RQC15: every clause of the original cusp-closure producer.}
\label{claim:RQC15}
For the same finite pair and any original \(\mathrm{IsReg}\), if \(\phi\)
in RQC13 satisfies the literal \(\mathrm{ZCusp}(V,\phi)\), every
continuous rational invariant representative \(\psi\) of a class
in its closed orbit span satisfies \(\mathrm{ZCusp}(V,\psi)\).
This is \texttt{ZClosed.zcusp\_of\_mem\_closure}.

\emph{Proof.}
Fix arbitrary \(\xi\) satisfying \(\mathrm{IsoVec}\), arbitrary
Haar measure \(\eta\) on \(ZI(V,\xi)\), and arbitrary original
rational fundamental domain \(\mathcal D\).
Take \(\mathcal E\subset C\) from RQC14 and push
\(\eta|_{\mathcal E}\) into \(G\) by the subgroup inclusion.
The resulting finite measure is supported on compact \(C\).
The hypothesis gives \(A_m\phi=0\) everywhere.  RQC13 gives
\(A_m\psi=0\) everywhere.  RQC14 transfers its integral from
\(\mathcal E\) to \(\mathcal D\), giving
\[
 \int_{\mathcal D}\psi(zg)\,d\eta(z)=0\qquad(g\in G).
 \tag{RQC.21}
\]
The three choices were arbitrary, exactly as in \(\mathrm{ZCusp}\).
If no nonzero isotropic vector exists, this condition is vacuous,
exactly as in the source definition.
% END RQC15

% BEGIN RQC16
\paragraph{RQC16: original finite level and genuine cusp for supplied representatives.}
\label{claim:RQC16}
Apply HCU18 with the regular representation RQC6.
Its original Hilbert map \(\Phi\) has RQC9--10's smoothness and
all-word identities.  If its seed \(\phi\) is \(\mathrm{ZCusp}\),
every supplied continuous rational invariant representative
\(\psi_y\) of \(\Phi(y)\) is \(\mathrm{ZCusp}\).
All original finite eigenlaws hold pointwise.  When those laws
include the original open finite level, \(\psi_y\) is
\(\mathrm{IsAutFun}\), and CTC12 makes it genuinely cuspidal.

\emph{Proof.}
HCU18 puts \(\Phi(y)\) in the closed span of the full adelic orbit
of the same seed.  Apply RQC15.  For each fixed finite \(k\), the
Hilbert eigenlaw says that \(\psi_y(\cdot k)\) and
\(\chi(k)\psi_y\) represent the same class.  HCU16's uniqueness of
continuous representatives makes them equal everywhere: their
continuous difference is globally almost everywhere zero, and
Haar measure has full support.  Rational invariance and continuity
are supplied, so the trivial-character restriction to the fixed
open finite level gives every field of \(\mathrm{IsAutFun}\).
CTC12 converts the exact transvection integrals to genuine cusp
integrals.  RQC10 gives the same conclusions for every derived tuple
for which a continuous representative is supplied, with the same
level.  Existence of these representatives is not inferred.
% END RQC16

\subsection{The original Gelfand operators act by scalars}

We reconstruct the scalar producer itself, using the actual smooth
extremal vector supplied by HCU12.  No theorem on infinitesimal
characters is assumed.  On the smooth subspace put
\(D(Z)=d\sigma_{\mathbb C}(Z)\), write \(E_{ab}\) for the usual
matrix units, and set
\[
 Z_m=\sum_{i_0,\ldots,i_{m-1}}
 D(E_{i_0i_1})D(E_{i_1i_2})\cdots D(E_{i_{m-1}i_0})
 \quad(m\ge1).
 \tag{RQC.22}
\]
Products retain their written composition order, with the rightmost
factor acting first.  Thus \(Z_{k+1}v\) is exactly
\(\mathrm{gelfRep}(H,\sigma,k,v)\), with no extra factor or sign.

% BEGIN RQC17
\paragraph{RQC17: noncommutative trace and change of matrix frame.}
\label{claim:RQC17}
Let \(\mathcal E=(D(E_{ab}))_{a,b}\) be a matrix in the
endomorphism algebra of smooth vectors.  Then
\(Z_m=\operatorname{tr}(\mathcal E^m)\).
For every \(p\in\mathrm{GL}_n(\mathbb C)\), replacing each
\(D(E_{ab})\) by \(D(pE_{ab}p^{-1})\) leaves \(Z_m\) unchanged.

\emph{Proof.}
Expand matrix multiplication without commuting entries.
The entry \((\mathcal E^m)_{ab}\) is the sum of ordered products
along length-\(m\) paths from \(a\) to \(b\).  Summing over
\(a=b\) gives (RQC.22).
Complex linearity from RQC8 gives, with a distinct name for the
conjugated operator matrix,
\[
 \mathcal E_p:=(D(pE_{ab}p^{-1}))_{a,b}
    =p^{\mathsf T}\mathcal E(p^{-1})^{\mathsf T}.
 \tag{RQC.23}
\]
The scalar entries of the outside matrices commute with every
operator entry.  Hence
\(\mathcal E_p^m=p^{\mathsf T}\mathcal E^m(p^{-1})^{\mathsf T}\)
by neighboring cancellation.
For a scalar matrix \(A\) and an operator-entry matrix \(B\),
summing diagonal entries gives
\(\operatorname{tr}(AB)=\operatorname{tr}(BA)\).
Use this once to cancel the outside scalar matrices.
The proof does not use cyclic invariance of trace for two
arbitrary noncommuting operator-entry matrices.
% END RQC17

% BEGIN RQC18
\paragraph{RQC18: scalar action on the original highest vector.}
\label{claim:RQC18}
Suppose a nonzero smooth \(v_0\) satisfies, in the original
invertible frame \(x_0\),
\[
 D(x_0E_{ab}x_0^{-1})v_0=0\ (a<b),\qquad
 D(x_0E_{aa}x_0^{-1})v_0=w_av_0.
 \tag{RQC.24}
\]
For every \(m\ge1\) there is \(c_m\) with \(Z_mv_0=c_mv_0\).
A finite integer-coefficient polynomial in the \(w_a\) supplies
such a scalar.

\emph{Proof.}
RQC17 permits work in the frame conjugated by \(x_0\).
In this paragraph include that conjugation in the matrix-unit
notation.  RQC8 and matrix-unit multiplication give
\[
 [D(E_{ab}),D(E_{cd})]
   =\delta_{bc}D(E_{ad})-\delta_{da}D(E_{cb}).
 \tag{RQC.25}
\]
Order the generators with strictly lower units first, diagonal
units second, and strictly upper units last; fix any order within
each list.  Every word is a linear combination of words in this
order.  To prove this, swap an adjacent inversion by
\(AB=BA+[A,B]\).  The swapped word has fewer inversions at the
same length, and (RQC.25) gives shorter commutator terms.
Induction on length and then inversion count terminates.
Only this spanning statement is needed; no independence of the
ordered words is assumed.

Assign weight \(e_a-e_b\) to \(D(E_{ab})\).
Every cycle in (RQC.22) has weight zero.  Each swap and every
nonzero commutator term preserve the total weight, so every
resulting ordered word has weight zero.
If such a word has an upper generator, the rightmost upper
generator kills \(v_0\).
If it has no upper generator but has a lower generator, its weight
cannot be zero: the functional
\(e_a-e_b\mapsto b-a\) is strictly negative on all lower units
and zero on the diagonal ones.
Thus the only surviving words consist entirely of diagonal
generators.  They act by products of the \(w_a\).
Their integer coefficients, introduced by (RQC.25), give the
claimed polynomial.  For \(n=0\), the cycle sum is empty and
\(c_m=0\).  The source endpoint asks for a scalar, so a particular
closed formula for this polynomial is unnecessary.
% END RQC18

% BEGIN RQC19
\paragraph{RQC19: symmetry and covariance of the exact cyclic operators.}
\label{claim:RQC19}
For smooth \(v,w\) and \(g\in U(H)\),
\[
 \langle Z_mv,w\rangle=\langle v,Z_mw\rangle,\qquad
 \sigma(g)Z_mv=Z_m\sigma(g)v.
 \tag{RQC.26}
\]
In particular, \(c_m\) in RQC18 is real.

\emph{Proof.}
RQC8 gives
\[
 \langle D(E_{ab})v,w\rangle
  =\langle v,D(H^{-1}E_{ba}H)w\rangle.
 \tag{RQC.27}
\]
Move the factors of each product in (RQC.22) successively into
the second argument.  Their order reverses and their indices
transpose.  Reindexing reversed closed paths produces the cyclic
trace formed from \(D(H^{-1}E_{ab}H)\).
RQC17 with \(p=H^{-1}\) makes it exactly \(Z_m\).
RQC7 permits all these successive pairings on smooth vectors,
proving symmetry.

Covariance in RQC8 moves \(\sigma(g)\) through the factors,
replacing each matrix unit by \(gE_{ab}g^{-1}\).
RQC17 makes the sum unchanged.
Finally, symmetry at \(v=w=v_0\) gives
\(\overline{c_m}\|v_0\|^2=c_m\|v_0\|^2\);
nonvanishing of \(v_0\) gives reality.
These are identities on the smooth domain, without claims of
boundedness or of a chosen self-adjoint extension.
% END RQC19

% BEGIN RQC20
\paragraph{RQC20: the full original Gelfand scalar producer.}
\label{claim:RQC20}
Assume \(\sigma\) is irreducible and unitary and that the actual
\(\mathrm{OrthFr}(H,q,x_0)\), nonzero smooth \(v_0\) and integral
weights \(w\) satisfy (RQC.24).
For every \(k\ge0\), one scalar \(c_{k+1}\) satisfies
\[
 \mathrm{gelfRep}(H,\sigma,k,v)=c_{k+1}v
 \quad\text{for every original smooth }v.
 \tag{RQC.28}
\]
This is every quantifier of \texttt{GelfScalar.gelfRep\_scalar}.

\emph{Proof.}
Taking determinants in \(x_0^*Hx_0=J\), with \(J\) the original
diagonal sign matrix, shows \(\det H\ne0\) and \(\det x_0\ne0\).
Then \(H=(x_0^{-1})^*Jx_0^{-1}\) is Hermitian.
RQC7--19 therefore apply.
The closed span of \(\{\sigma(g)v_0:g\in U(H)\}\) is a nonzero
closed invariant subspace, hence the whole Hilbert space by
irreducibility.
For any smooth \(v\), symmetry, covariance and the scalar action
on \(v_0\) give
\[
 \langle(Z_{k+1}-c_{k+1})v,\sigma(g)v_0\rangle=0
 \quad(g\in U(H)).
 \tag{RQC.29}
\]
The realness of \(c_{k+1}\) ensures cancellation with the stated
inner-product convention.
Density forces the fixed vector
\((Z_{k+1}-c_{k+1})v\) to be zero.
No unbounded operator has been extended by continuity: its
already defined value on a single smooth vector is paired with a
dense set.  Formula (RQC.22) identifies this operator with the
literal source \(\mathrm{gelfRep}\), proving the endpoint.
% END RQC20

% BEGIN RQC21
\paragraph{RQC21: exact pointwise central finiteness once derivatives are supplied.}
\label{claim:RQC21}
For the original model with \(\mathrm{HolShape}(M)\),
\(\mathrm{WtIs}(M,w)\) and a supplied original
\(\mathrm{IsExtVec}(M,v_0)\) tuple, as in HCU12, suppose
\[
 T:\prod_\tau M.\mathrm{sm}_\tau\longrightarrow(G\to\mathbb C)
\]
is complex multilinear and has the pointwise identities
\[
 R\mathrm{der}(\tau,X)(T(y))
  =T\bigl(y[\tau\leftarrow d(M.\mathrm{rep}_\tau)(X)y_\tau]\bigr)
 \quad(X\in\mathrm{lieAt}(V,\tau)).
 \tag{RQC.30}
\]
For every \(\tau,k\), one scalar \(c_{\tau,k}\) then satisfies
\[
 \mathrm{gelf}(\theta,V,\tau,k)(T(y))=c_{\tau,k}T(y)
 \quad\text{for all }y.
 \tag{RQC.31}
\]
The span \(\mathbb C\,T(y)\), of dimension at most one, witnesses
the original \(\mathrm{IsAutFunSm.zfin}\) field.

\emph{Proof.}
HCU12 gives the actual frame, nonzero smooth extremal vector,
Borel equations and irreducible unitary model at each place.
RQC20 therefore gives one scalar on every member of the original
\(M.\mathrm{sm}_\tau\); RQC7 identifies its span definition with
smooth vectors.
Decompose a complex direction as in (RQC.11), apply (RQC.30)
to its two real parts, and use multilinearity in that coordinate.
This is the literal \(R\mathrm{derC}\) identity.
Induction gives the identity for every ordered complex word.
Summing over the closed matrix-unit paths of length \(k+1\)
gives \(\mathrm{gelfRep}\) in the model coordinate, and RQC20
gives (RQC.31).  This also proves the exact
\texttt{rderC\_value}, \texttt{wordC\_value} and
\texttt{gelf\_value} consumers.

For a scalar multiple \(aT(y)\), write it as
\(T(y[\tau\leftarrow ay_\tau])\) and apply the same identity.
Thus every original Gelfand operator preserves \(\mathbb C\,T(y)\).
If the place set is empty, the stability condition is vacuous
and the same span works.
This is precisely \texttt{BackForm.zfin\_value}.
The existence of \(T\) with (RQC.30) is still an input.
% END RQC21

% BEGIN RQC22
\paragraph{RQC22: the exact remaining boundary after Hilbert realization.}
\label{claim:RQC22}
For the original supplied bounded \(\mathrm{HolForm}(\theta,M,w,\ldots,\phi)\),
\(\mathrm{FinPair}\), \(\mathrm{IsInfFam}(\theta)\),
\(\mathrm{HolShape}(M)\), \(\mathrm{WtIs}(M,w)\) and
\(\mathrm{IsExtVec}(M,v_0)\),
the HCU18 Hilbert realization needs no additional
\(\mathrm{IsReg}\) assumption.
It has all the original smooth-vector and derivative-word
properties.  Once continuous smooth representatives and their
pointwise derivative identities are supplied, their original
finite laws and central finiteness follow from RQC16 and RQC21.
If the same seed also satisfies \(\mathrm{ZCusp}(V,\phi)\),
their genuine cuspidality follows as well.

\emph{Proof.}
Construct \(\rho\) by RQC6 and insert it into HCU16--18.
The result uses the same quotient measure (RQC.1), seed class,
model and \(\mathrm{atTau}\).  RQC9--10 apply to this output.
Its closed-orbit containment and finite eigenlaws are the ones in
HCU18, so RQC16 supplies the finite-level conclusion and,
under the stated seed \(\mathrm{ZCusp}\) hypothesis, the cusp
conclusion when representatives are given.  Under pointwise derivative
identities, HCU12 and RQC20--21 supply central finiteness with
the exact original normalization.
Thus every conclusion belongs to this same map, without
replacing it by an abstract isomorphic model.

The still-unproved production includes the target finite pair,
initial boundedness, and actual smooth representatives with
Sobolev estimates giving pointwise differentiation, continuous
evaluation and uniform moderate growth.  These remain in the
closure record.  The full \(\mathrm{IsMultMap}\) and full theta
extension are not asserted here.  In particular the post-Hilbert
scalar argument is not used to justify earlier boundedness;
the realized Hilbert map does not supply one of its own inputs.
% END RQC22
\end{document}
